%% kapa: exam=1 round=28 year=2017 period=2 form=A subject=accounting subject_ko=회계학 seq=40 number=80
%% source: Q-Net 감정평가사 기출문제 (cst00309) artlSeq=5199603 file=extracted/1008/5199603/2017년 제28회 감정평가사 1차시험 문제지/2교시 A형.pdf
%% conversion: kordoc 4.24.51 -> tools/gen1.py (auto); review=proofread-sonnet; flags=none
\begin{question}{80}
(주)감평은 20x1년 1월 1일에 설립된 회사이다. 20x1년도 1월 및 2월의 원가자료는 다음과 같다.
\par{\small\centering\begin{tabular}{|l|r|r|}
\hline
\multicolumn{1}{|c|}{구분} & \multicolumn{1}{c|}{1월} & \multicolumn{1}{c|}{2월} \\ \hline
최대생산가능량 & 1,000단위 & 1,200단위 \\ \hline
생산량 & 800단위 & 1,000단위 \\ \hline
판매량 & 500단위 & 1,100단위 \\ \hline
변동제조원가(총액) & ￦40,000 & ￦50,000 \\ \hline
고정제조간접원가(총액) & ￦20,000 & ￦30,000 \\ \hline
변동판매관리비(총액) & ￦1,500 & ￦5,500 \\ \hline
고정판매관리비(총액) & ￦2,000 & ￦2,000 \\ \hline
\end{tabular}\par}
(주)감평은 실제원가계산을 적용하고 있으며, 원가흐름가정은 선입선출법이다. 20x1년 2월의 전부원가계산에 의한 영업이익이 ￦10,000이면, 2월의 변동원가계산에 의한 영업이익은? (단, 기초 및 기말 재공품재고는 없다.) \par
\begin{choices}
  \item ￦10,500
  \item ￦11,000
  \item ￦11,500
  \item ￦12,000
  \item ￦12,500
\end{choices}
\end{question}
\answer{3}
